\section{Breadth-First Search}

Breadth-First Search (BFS) is a graph-search algorithm that explores nodes
level by level outward from a start node.

BFS is designed for unweighted graphs, or equivalently for graphs in which
every edge has the same cost.

Its key property is that the first path it finds to a node uses the minimum
possible number of edges from the start.

Therefore, BFS minimizes

\[
\text{hop count},
\]

rather than weighted path cost.

\subsection{Queue-Based Exploration}

BFS uses a first-in, first-out queue.

The start node is inserted into the queue first.

The algorithm repeatedly:

\begin{enumerate}
    \item removes the node at the front of the queue,
    \item examines all of its unvisited neighbors,
    \item marks newly discovered neighbors as visited,
    \item records the current node as their parent,
    \item inserts those neighbors at the back of the queue.
\end{enumerate}

Because nodes are processed in the same order in which they are discovered,
the search expands the graph one level at a time.

Conceptually, the traversal proceeds as

\[
\text{level }0
\rightarrow
\text{level }1
\rightarrow
\text{level }2
\rightarrow
\cdots
\]

where level $k$ contains nodes that can be reached from the start using
exactly $k$ edges.

\subsection{Visited Nodes}

A visited structure is used to prevent nodes from being inserted into the
queue multiple times.

In the current project, graph nodes are indexed by dense integer IDs,

\[
0,1,\ldots,|V|-1,
\]

so a boolean array can be used:

\[
\texttt{visited}[i].
\]

A node is marked visited when it is inserted into the queue.

This prevents two different already-expanded nodes from inserting the same
neighbor into the queue more than once.

\subsection{Shortest Path in an Unweighted Graph}

Suppose the goal node is first discovered at level $k$.

Every node in levels

\[
0,1,\ldots,k-1
\]

has already been explored before any node at level $k+1$ is considered.

Therefore, no path containing fewer than $k$ edges can exist.

Hence BFS returns a path with minimum hop count.

If all graph edges have identical cost $c$, minimizing hop count also
minimizes total path cost because

\[
\text{cost}(P)
=
c \cdot
\text{number of edges in }P.
\]

\subsection{BFS on a Weighted Visibility Graph}

The visibility graph used in this project is weighted.

Each edge weight is its Euclidean length:

\[
w(u,v)
=
\|u-v\|.
\]

BFS ignores these edge weights during search.

It therefore optimizes

\[
\text{number of visibility edges}
\]

rather than

\[
\text{total Euclidean path length}.
\]

This means a BFS path can use fewer graph edges while still being
geometrically longer than the path found by Dijkstra's algorithm or A*.

For example, BFS may prefer a direct one-edge route of length

\[
10
\]

over a four-edge route whose total length is

\[
4.
\]

BFS chooses the first path because

\[
1 < 4,
\]

while Dijkstra chooses the second because

\[
4 < 10.
\]

This difference makes BFS useful in the project as a comparison algorithm,
even though it is not the correct weighted shortest-path planner for the
visibility graph.

\subsection{Path Reconstruction}

Like Dijkstra's algorithm and A*, BFS stores a parent for each newly
discovered node.

If node $v$ is first discovered while expanding node $u$, then

\[
\text{parent}(v)=u.
\]

Once the goal is reached, parent links are followed backward:

\[
G
\rightarrow
\text{parent}(G)
\rightarrow
\cdots
\rightarrow
S.
\]

This produces the path in reverse order.

Reversing the sequence gives

\[
S
\rightarrow
\cdots
\rightarrow
G.
\]

\subsection{Reporting Geometric Path Length}

Although BFS ignores edge weights during search, the implementation still
reports the total geometric length of the final path.

If the BFS path contains edges

\[
e_1,e_2,\ldots,e_k,
\]

then the reported geometric distance is

\[
L_{\text{BFS}}
=
\sum_{i=1}^{k} w(e_i).
\]

This value is measured after the path has been selected.

It does not influence the BFS search order.

Reporting geometric path length allows BFS results to be compared directly
with Dijkstra and A* in the interactive planner.

\subsection{Undirected Visibility Edges}

Visibility is symmetric.

If node $u$ is visible from node $v$, then node $v$ is also visible from
node $u$.

The visibility graph therefore represents an undirected graph.

Each undirected edge is stored only once, but the BFS adjacency list inserts
both directions:

\[
u \rightarrow v
\]

and

\[
v \rightarrow u.
\]

This allows BFS to traverse the visibility graph correctly regardless of
which endpoint appears first in the stored edge.

\subsection{Expanded Nodes}

For consistency with Dijkstra and A*, this project counts a node as expanded
when it has been removed from the queue and its neighbors are about to be
examined.

The goal node is not counted as expanded if the search terminates immediately
when the goal is removed from the queue.

This shared definition allows expanded-node counts to be compared across all
three planners.

\subsection{Complexity}

Using an adjacency-list representation, every vertex is inserted into and
removed from the queue at most once.

Every undirected edge is examined a constant number of times.

Therefore, the time complexity is

\[
O(|V|+|E|).
\]

The visited array, parent array, adjacency list, and queue require

\[
O(|V|+|E|)
\]

space overall.

If only the additional search state beyond the graph representation is
counted, the BFS working memory is

\[
O(|V|).
\]

\subsection{Comparison with Dijkstra and A*}

The three current graph-search planners differ primarily in how they choose
which node to explore.

BFS processes nodes according to graph depth and minimizes

\[
\text{hop count}.
\]

Dijkstra selects the node with the smallest known weighted cost

\[
g(n)
\]

and minimizes total path weight.

A* selects nodes using

\[
f(n)=g(n)+h(n),
\]

where $h(n)$ estimates the remaining cost to the goal.

For the current Euclidean visibility graph:

\begin{itemize}
    \item BFS is not guaranteed to produce the geometrically shortest path.
    \item Dijkstra and A* should produce the same optimal geometric path cost.
    \item A* may expand fewer nodes than Dijkstra because of its heuristic.
\end{itemize}

This makes BFS, Dijkstra, and A* useful together for demonstrating the effect
of different graph-search objectives and search strategies.